\noindent All figures in this section are derived from the curated
\textsc{AniSpeakArchive} (ASA) corpus, a domain-specific literature archive and
discovery platform for animal communication research.
\HL{placeholder comment cite AniSpeakArchive here}
The methodology, terminology, category definitions, and all descriptive statistics
reported in the plots below are derived from the AniSpeakArchive corpus and its
annotation framework.
\HL{Need to discuss the methodology in more detail here.}
 % New plots by H
\paragraph{Publication Growth and the 2000 Contrast Point.}

\figref{fig:publication_trends} shows that the traditional literature provides the
long historical backbone of the field, whereas computational and hybrid work remains
limited until the early 2000s and then expands much more visibly in the following
years. The most pronounced acceleration appears after roughly 2015, consistent with
the broader uptake of deep learning, large-scale audio processing pipelines, and more
widely available open-source machine learning tools in bioacoustics and related areas.
We therefore use 2000 as a coarse corpus-level contrast point for the comparative
analyses below, while noting that the strongest recent inflection in
computational-oriented growth occurs in the mid-2010s. The terminal bin (2025+) is
partial because metadata consolidation and indexing typically lag the nominal
publication year.

\begin{figure}[th]
  \centering
  \includegraphics[width=0.95\linewidth]{figures_trends_new/publication_trends_two_groups.pdf}
  \caption{Annual publication counts in \textsc{AniSpeakArchive} (1950--2026),
    stratified into traditional and computational-oriented research traditions.}
  \label{fig:publication_trends}
\end{figure}

%% ── Authorship ───────────────────────────────────────────────────────────────
\paragraph{Authorship Patterns.}

\figref{fig:top_authors} indicates that the traditional and computational-oriented
literatures are not simply two views of the same author community. The traditional
side is broader and more historically accumulated, whereas the computational-oriented
side is anchored by a smaller set of recurring contributors, suggesting a more
specialized methodological community at the intersection of animal behavior,
bioacoustics, signal processing, and machine learning.

\begin{figure}[th]
  \centering
  \includegraphics[width=0.85\linewidth]{figures_trends_new/top_authors_dual.pdf}
  \caption{Most prolific authors in \textsc{AniSpeakArchive}, shown separately for
    traditional and computational-oriented research traditions.}
  \label{fig:top_authors}
\end{figure}

%% ── Venues and institutions ──────────────────────────────────────────────────
\paragraph{Publication Venues and Institutional Affiliations.}

\figref{fig:top_venues} and \figref{fig:top_institutions} show that both traditions
share some core publication infrastructure, but they are not distributed identically
across venues or institutions. The logarithmic scaling in the venue plot makes clear
that a small number of outlets dominate the archive, with the
\textit{Journal of the Acoustical Society of America} occupying a particularly central
position. At the same time, the institutional profile suggests that computationally
oriented work is concentrated in a narrower subset of research groups, again
consistent with a more specialized technical subcommunity.

\begin{figure}[th]
  \centering
  \begin{subfigure}{\linewidth}
    \centering
    \includegraphics[width=0.95\linewidth]{figures_trends_new/top_venues_dual.pdf}
    \caption{Most frequent publication venues (log scale).}
    \label{fig:top_venues}
  \end{subfigure}
  \par\vspace{0.9em}
  \begin{subfigure}{\linewidth}
    \centering
    \includegraphics[width=0.95\linewidth]{figures_trends_new/top_institutions_dual.pdf}
    \caption{Most productive institutional affiliations.}
    \label{fig:top_institutions}
  \end{subfigure}
  \caption{Publication venues (top) and institutional affiliations (bottom) in
    \textsc{AniSpeakArchive}, each stratified by research tradition.}
  \label{fig:venues_institutions}
\end{figure}

%% ── Taxa and features ────────────────────────────────────────────────────────
\paragraph{Taxa and Communication Features.}

\figref{fig:top_species} shows substantial overlap in organismal focus across the two
traditions, with cetaceans and bats occupying especially prominent positions in both.
This likely reflects the strong role of acoustically rich and methodologically tractable
systems in the archive. Passerine songbirds appear relatively more prominent in the
traditional literature, consistent with their longstanding importance in behavioral and
comparative studies of vocal communication. \figref{fig:top_features} shows that the
two traditions also differ in analytic emphasis: traditional work is more closely tied
to established biological descriptors of signaling behavior, whereas the
computational-oriented literature places greater weight on features that are amenable
to automated extraction, representation, and classification.

\begin{figure}[th]
  \centering
  \includegraphics[width=0.85\linewidth]{figures_trends_new/top_species_dual.pdf}
  \caption{Most frequently studied taxa by research tradition in
    \textsc{AniSpeakArchive}. Labels reflect normalised common names; the catch-all
    \emph{Other} category is excluded.}
  \label{fig:top_species}
\end{figure}

\begin{figure}[th]
  \centering
  \includegraphics[width=0.85\linewidth]{figures_trends_new/top_features_dual_pct.pdf}
  \caption{Most frequently analysed acoustic and linguistic features by research
    tradition in \textsc{AniSpeakArchive}, expressed as shares of all feature
    mentions within each tradition.}
  \label{fig:top_features}
\end{figure}

%% ── Corpus composition ───────────────────────────────────────────────────────
\paragraph{Overall Corpus Composition.}

\figref{fig:domains_donut}, \figref{fig:contexts_donut}, and
\figref{fig:modalities_donut} summarize the corpus at the global level, irrespective
of tradition or period. Together, these plots show that the archive is heavily shaped
by acoustically centered research, both in modality and in application focus. The
context distribution further indicates that field-based and naturalistic studies remain
the dominant empirical setting, even though computational analyses have become more
visible. This overall composition is important for interpreting later comparisons,
because methodological trends in the archive are still grounded primarily in vocal and
audio-based animal communication research.

\begin{figure}[th]
  \centering
  \includegraphics[width=0.72\linewidth]{figures_trends_new/domains_donut.pdf}
  \caption{Distribution of application domains across the full
    \textsc{AniSpeakArchive} corpus.}
  \label{fig:domains_donut}
\end{figure}

\begin{figure}[th]
  \centering
  \includegraphics[width=0.72\linewidth]{figures_trends_new/contexts_donut.pdf}
  \caption{Distribution of research settings across the full
    \textsc{AniSpeakArchive} corpus.}
  \label{fig:contexts_donut}
\end{figure}

\begin{figure}[th]
  \centering
  \includegraphics[width=0.72\linewidth]{figures_trends_new/modalities_donut.pdf}
  \caption{Distribution of communication modalities across the full
    \textsc{AniSpeakArchive} corpus.}
  \label{fig:modalities_donut}
\end{figure}

%% ── Lexical shifts ───────────────────────────────────────────────────────────
\paragraph{Lexical Shifts Before and After 2000.}

To complement the count-based summaries above, we examine how abstract-level
vocabulary changes across the 2000 contrast point using TF-IDF scoring and word
clouds. \figref{fig:tfidf_top_terms} highlights the dominant terms within each period,
revealing a shift from earlier biologically oriented vocabularies toward a more
methodologically explicit lexicon in the later period. \figref{fig:tfidf_delta}
makes this transition more explicit by identifying terms whose relative prominence
increases or decreases most strongly after 2000. \figref{fig:wordcloud_prepost}
provides a complementary qualitative view, making it easier to see how the later
literature becomes more saturated with technical and computational terminology.

\begin{figure}[th]
  \centering
  \includegraphics[width=0.88\linewidth]{figures_trends_new/tfidf_top_terms_2000.pdf}
  \caption{Dominant abstract terms before and after 2000, ranked by mean TF-IDF
    within each period.}
  \label{fig:tfidf_top_terms}
\end{figure}

\begin{figure}[th]
  \centering
  \includegraphics[width=0.72\linewidth]{figures_trends_new/tfidf_delta_2000.pdf}
  \caption{Change in mean TF-IDF across the 2000 contrast point (post-2000 minus
    pre-2000). Positive values indicate terms that gained relative prominence after
    2000; negative values indicate decline.}
  \label{fig:tfidf_delta}
\end{figure}

\begin{figure}[th]
  \centering
  \begin{subfigure}{0.49\linewidth}
    \centering
    \includegraphics[width=\linewidth]{figures_trends_new/wordcloud_pre2000.pdf}
    \caption{Pre-2000}
    \label{fig:wordcloud_pre2000}
  \end{subfigure}\hfill
  \begin{subfigure}{0.49\linewidth}
    \centering
    \includegraphics[width=\linewidth]{figures_trends_new/wordcloud_post2000.pdf}
    \caption{Post-2000}
    \label{fig:wordcloud_post2000}
  \end{subfigure}
  \caption{Abstract-term word clouds before and after 2000 in \textsc{AniSpeakArchive}.
    Word size reflects relative prominence within each period.}
  \label{fig:wordcloud_prepost}
\end{figure}
